\section{Task 4: Comparison and Interpretation}

\subsection{ca-GrQc: Real vs ER vs BA}
The real ca-GrQc network has substantially greater degree heterogeneity than ER: their degree variances are 62.690 and 5.157, respectively (Table~\ref{tab:global-distances}). BA is closer on this measure, with variance 74.714, but its maximum degree of 266 exceeds the real network's 81 (Figure~\ref{fig:degree-grqc}). Thus, preferential attachment provides a more suitable description of unequal connectivity than independent random edges in this comparison, without reproducing the full observed degree distribution or establishing how the real network formed.

Neither model captures the strong local clustering. Mean clustering is 0.529636 in the real network, compared with 0.000921 for ER and 0.009484 for BA (Table~\ref{tab:distribution-summary}; Figure~\ref{fig:clustering-grqc}). The high real-network clustering is consistent with its construction: coauthors of a multi-author paper form a clique. Neither reference model explicitly includes this group-based construction.

The models also underrepresent fragmentation and LCC distances. The real graph has 355 components, versus 19 for ER and one for BA (Table~\ref{tab:global-connectivity}; Figure~\ref{fig:components-grqc}). Mean eccentricity is 11.675 in the real LCC, versus 8.047 and 5.686, respectively (Figure~\ref{fig:ecc-grqc}). These differences indicate that matching node counts and approximate edge counts does not reproduce component structure or distances within the largest component. BA captures degree heterogeneity more closely, but neither model jointly reproduces clustering, fragmentation, and eccentricity.

\subsection{Western States Power Grid: Real vs ER vs BA}
ER is closer to the Power Grid's degree variability: its variance is 2.465, compared with 3.209 for the real graph and 16.740 for BA (Table~\ref{tab:global-distances}). Maximum degrees are 11, 19, and 206 for ER, real, and BA, respectively (Figure~\ref{fig:degree-power}). BA therefore introduces much stronger hub dominance than observed, whereas ER better approximates the relatively narrow degree distribution.

This agreement does not extend to connectivity or distances. The real grid is connected, while ER has 387 components, including 345 isolates; BA is connected (Table~\ref{tab:global-connectivity}; Figure~\ref{fig:components-power}). Nevertheless, both models have much smaller mean eccentricities: 14.924 for ER and 14.390 for BA, versus 34.541 for the real grid (Figure~\ref{fig:ecc-power}). ER distances describe only its LCC. The real grid's longer distances are consistent with geographically constrained infrastructure, although these measurements do not establish that explanation causally.

Mean clustering is 0.080104 in the real grid, 0.000502 in ER, and zero in BA (Figure~\ref{fig:clustering-power}). BA's zero clustering follows from its $m=1$ tree construction. Its 25.08\% edge deficit (Table~\ref{tab:model-parameters}) also confounds structural comparisons, so differences cannot be attributed solely to attachment rules. ER better approximates degree variability and BA preserves connectedness, but neither reproduces the grid's clustering and eccentricity. Across both datasets, conclusions concern single seeded realizations with unequal realized edge counts, not ensemble-level model performance.
\clearpage
